\documentclass[12pt,a4paper]{article}

\usepackage[margin=2.4cm]{geometry}
\usepackage{fontspec}
\usepackage{xeCJK}
\usepackage{graphicx}
\usepackage{xcolor}
\usepackage{booktabs}
\usepackage{longtable}
\usepackage{array}
\usepackage{enumitem}
\usepackage{hyperref}
\usepackage{titlesec}
\usepackage{fancyhdr}
\usepackage{setspace}

\setmainfont{Times New Roman}
\setsansfont{Arial}
\setmonofont{Menlo}
\setCJKmainfont{PingFang TC}
\setCJKsansfont{PingFang TC}
\setCJKmonofont{PingFang TC}

\definecolor{linquestGreen}{HTML}{3F5F4A}
\definecolor{linquestWarm}{HTML}{F4EFE3}
\definecolor{linquestAccent}{HTML}{9A5A45}
\definecolor{linquestText}{HTML}{243229}

\hypersetup{
  colorlinks=true,
  linkcolor=linquestGreen,
  urlcolor=linquestAccent,
  citecolor=linquestGreen
}

\pagestyle{fancy}
\fancyhf{}
\lhead{Linquest 期末報告}
\rhead{\thepage}
\renewcommand{\headrulewidth}{0.4pt}

\titleformat{\section}{\Large\bfseries\color{linquestGreen}}{\thesection}{0.8em}{}
\titleformat{\subsection}{\large\bfseries\color{linquestGreen}}{\thesubsection}{0.8em}{}
\titleformat{\subsubsection}{\normalsize\bfseries\color{linquestText}}{\thesubsubsection}{0.8em}{}

\setstretch{1.2}
\setlist[itemize]{leftmargin=2em, itemsep=0.25em, topsep=0.3em}
\setlist[enumerate]{leftmargin=2em, itemsep=0.25em, topsep=0.3em}
\renewcommand{\arraystretch}{1.25}
\setlength{\tabcolsep}{3.5pt}
\emergencystretch=3em

\newcolumntype{P}[1]{>{\raggedright\arraybackslash}p{#1}}

\begin{document}

\begin{titlepage}
  \pagecolor{linquestWarm}
  \centering
  \vspace*{2.4cm}
  {\Huge\bfseries\color{linquestGreen} Linquest 期末報告書\par}
  \vspace{0.8cm}
  {\Large 語言學習、班級領地競賽與個人化練習系統\par}
  \vspace{2cm}
  \begin{tabular}{rl}
    專案名稱： & Linquest \\
    組員名單： & 楊立宇 --- 題材發想、程式開發、系統整合 \\
              & 顏端廷 --- 題材延伸、程式開發、功能設計 \\
    提交日期： & 2026 年 5 月 31 日 \\
    開發平台： & Expo / React Native / TypeScript / Supabase \\
  \end{tabular}
  \vfill
  {\large 期中報告延伸版：由概念規劃進展至可執行主流程與後端資料模型\par}
  \nopagecolor
\end{titlepage}

\tableofcontents
\newpage

\section{摘要}

Linquest 是一款結合語言學習、班級合作競賽與個人自學路線的行動學習應用程式。期中階段的核心構想是把單字背誦從個人反覆練習，轉化為具有團隊目標與策略選擇的學習活動；期末階段則已進一步完成主要程式架構、學生端流程、教師後台、Supabase 資料表與測試基礎。

本專案的主要目的，是解決語言學習中常見的三個問題。第一，單字學習容易變成短期記憶，缺乏持續追蹤與錯題回流。第二，傳統測驗缺乏互動性，學生容易覺得重複刷題枯燥。第三，教師雖然可以指定單字範圍，卻不容易把全班學習狀態轉換成可觀察、可回饋的課堂活動。

Linquest 因此設計兩條主線。第一條是教師帶班的領地對戰模式：教師建立活動、匯入本次題庫、把學生分組，學生透過答題佔領六邊形地圖上的區塊，團隊依領地、財政與排行榜競爭。第二條是學生個人的 roadmap 自學模式：學生沿著關卡路線完成題組，系統記錄答題結果，並在同一輪練習中追問錯題，形成最低限度的個人化學習循環。

目標使用者包括國高中生、大學生、準備語言檢定的自學者，以及需要設計課堂互動活動的語言教師。相較於一般單字卡工具，Linquest 更強調班級合作、即時互動與教師活動管理；相較於固定課程型語言 app，Linquest 則保留教師自訂題庫與課堂情境整合的彈性。

\section{市場研究與正當性}

\subsection{競爭者分析}

\begin{longtable}{P{2.5cm}P{4.0cm}P{3.7cm}P{4.4cm}}
\toprule
產品 & 主要功能 & 優勢 & 限制與 Linquest 的改進方向 \\
\midrule
Quizlet & 單字卡、Learn、Test、Flashcards 等學習模式，支援使用者自建學習集。 & 上手容易，適合快速建立詞彙與測驗內容；題庫開放性高。 & 遊戲化主要停留在練習模式切換，較缺乏班級級別的長期合作目標。Linquest 以領地地圖、組別財政、特殊格對戰，把答題轉成團隊競賽與課堂事件。 \\
\midrule
Duolingo & 以短課程、連續學習 streak、遊戲化任務與提醒促進日常練習。 & 黏著度高，能用簡短任務降低開始學習的門檻。 & 課程內容較固定，不適合教師快速匯入某次段考或課堂指定單字。Linquest 允許教師為每輪活動建立自訂 CSV 題庫，讓遊戲活動直接對齊課堂內容。 \\
\bottomrule
\end{longtable}

Quizlet 的官方說明列出 Flashcards、Learn、Test 等模式，說明其優勢在於學習集與練習形式的彈性。Duolingo 官方則強調 game-like features、challenge、streak 與提醒，說明其產品重點是持續使用習慣。Linquest 參考兩者優勢，但把產品定位集中在「教師可控的班級活動」與「團隊領地競爭」，讓遊戲化不只是個人積分，而是能被教師安排到課堂中的學習任務。

\subsection{目標受眾定義}

本專案的主要使用者分為學生與教師兩類。

\begin{itemize}
  \item 學生：國高中生、大學生、語言檢定考生與自學者。他們需要反覆接觸單字，也需要比單純背誦更有動機的練習形式。
  \item 教師：需要帶班、派發單字範圍、檢視全班學習狀況的語言教師。教師可以把一份單字表轉為一輪活動，讓學生在活動期間以小組形式完成練習。
\end{itemize}

此應用程式對學習過程的幫助在於：學生不只是在完成個人題目，而是在為團隊地圖推進做出貢獻；教師不只是在發派測驗，而是能把題庫、分組、活動時間、排行榜與結算整合成一個完整教學活動。

\subsection{獨特性與創新點}

Linquest 的創新點可整理為三項。

\begin{enumerate}
  \item 將單字作答結果轉化為領地擴張、財政收入與團隊排名，使學習成果具備可視化的集體目標。
  \item 將教師自訂題庫納入活動生命週期，讓每一次課堂活動都能對應實際教材，而非只能使用固定官方內容。
  \item 在特殊格導入 1 對 1 即時對戰，使學生之間產生真實互動；但獲勝收益仍回到團隊，維持合作學習的核心。
\end{enumerate}

\section{應用程式功能與需求}

\subsection{核心功能}

\subsubsection{帳號與角色系統}

系統使用 Supabase Auth 實作 Email / password 登入、註冊與忘記密碼流程。註冊時可選擇 student 或 teacher 角色，登入後依角色進入不同主要介面。學生以 roadmap、territory 活動與個人頁面為主；教師則進入 console 管理班級與活動。

\subsubsection{roadmap 自學模式}

roadmap 是學生平時練習的模式。學生會看到關卡式路線，進入某一關後系統從 Google Sheets 題庫抓取該 level 的題目，並以四選一方式出題。每關以 15 題為基本單位，首輪正確率達 80\% 以上即可解鎖下一關；答錯題會被放回當輪 queue，直到錯題被再次完成。這讓 roadmap 雖然還不是完整 SRS，仍已具備錯題回流與進度解鎖的學習節奏。

\subsubsection{教師班級與活動管理}

教師後台可以建立班級、產生班級代碼、查看 roster、建立活動草稿、貼上或上傳 CSV 題庫、設定組數與活動結束時間。活動發布時，系統會依班級成員做隨機且人數平衡的分組，並產生對應 territory 地圖。

\subsubsection{領地對戰模式}

領地活動以六邊形地圖呈現。每組一開始有初始領地，學生只能從己方接壤區域往外挑戰。一般格透過答題佔領；倍率格提供更高價值，但搭配較高成本或難度；特殊格需透過 1 對 1 對戰取得。系統使用財政模型記錄攻擊成本、佔領獎勵、失敗懲罰與週期稅收，避免單純刷分造成遊戲失衡。

\subsubsection{即時 1 對 1 對戰}

特殊格會引導學生邀請其他組在線玩家進入 battle room。對戰使用 Supabase Realtime 與 presence 同步雙方狀態。雙方在同一組題目中答題，系統依分數判定勝者，勝方為其團隊取得特殊格收益。這是整個系統中唯一需要強同步的子流程，其餘地圖與排行榜更新則以一般查詢或 polling 處理。

\subsubsection{結算與輕量分析}

活動結束後，學生可查看結算與排行榜；教師可查看本次活動統計，例如常錯題、平均正確率與個別學生表現。活動資料採生命週期設計：教師可在結束後保留資料供檢視，也可主動刪除活動與其自訂題庫。

\subsection{使用者流程}

\subsubsection{學生流程}

\begin{enumerate}
  \item 學生註冊或登入。
  \item 若要參加課堂活動，輸入班級代碼加入班級。
  \item 進入 territory 活動列表，選擇一輪活動。
  \item 在地圖上查看己方領地、可挑戰格與排行榜。
  \item 點選可攻擊的格子，完成四選一題目。
  \item 若為特殊格，邀請其他組玩家進行 1 對 1 對戰。
  \item 活動結束後查看團隊排名、領地狀態與個人表現。
  \item 平時可進入 roadmap 自學模式，完成關卡並解鎖下一階段。
\end{enumerate}

\subsubsection{教師流程}

\begin{enumerate}
  \item 教師註冊或登入，進入 console。
  \item 建立班級，將班級代碼提供給學生。
  \item 建立新活動草稿，設定活動名稱、結束時間、組數與題庫。
  \item 貼上或上傳本輪 CSV 題庫，格式為中文提示、英文答案與選填詞性。
  \item 發布活動，系統自動分組並建立地圖。
  \item 活動期間查看排行榜、活動進度與學生參與情形。
  \item 活動結束後查看結算與輕量分析，必要時刪除活動。
\end{enumerate}

\subsection{使用的技術}

\begin{longtable}{P{4cm}P{10.5cm}}
\toprule
項目 & 使用技術與說明 \\
\midrule
前端框架 & Expo Router、React Native、React Native Web、TypeScript。單一 codebase 同時支援學生行動端與教師 web console。 \\
後端服務 & Supabase Auth、Postgres、Realtime、RLS、RPC。Auth 負責登入與角色；Postgres 儲存班級、活動、題庫、地圖、作答與對戰；Realtime 用於 1 對 1 battle 與 presence。 \\
題庫來源 & roadmap 題庫由 Google Sheets CSV runtime 讀取；教師活動題庫在建立活動草稿時寫入 Supabase custom bank。 \\
測試工具 & Jest、ts-jest、jest-expo、Testing Library for React Native。測試分為 unit tests 與需要 hosted Supabase test project 的 integration tests。 \\
主要資料表與模組 & users、classes、activities、question\_banks、questions、groups、hex\_tiles、attempts、battles、territory\_events、roadmap\_progress 等。 \\
\bottomrule
\end{longtable}

\section{UI/UX 設計}

\subsection{線框圖或模型}

期末階段已有多個主要畫面進入程式碼實作，包括：

\begin{itemize}
  \item 登入、註冊與忘記密碼畫面：\texttt{app/(auth)/}
  \item 學生首頁：\texttt{app/(app)/home/index.tsx}
  \item roadmap 路線、關卡與結果頁：\texttt{app/(app)/roadmap/}
  \item territory 活動列表、地圖、排行榜與結算頁：\texttt{app/(app)/territory/}
  \item 即時對戰房：\texttt{app/battle/[battleId].tsx}
  \item 教師 console 班級、活動與活動詳情頁：\texttt{app/(app)/console/}
\end{itemize}

下列圖為專案設計參考圖，用於確認整體視覺方向與介面氛圍。

\begin{figure}[h]
  \centering
  \includegraphics[width=0.78\linewidth]{docs/img_refs/37AC414D-2664-4104-95D6-E2B1A5700AF1.PNG}
  \caption{Linquest 設計參考圖之一：古典制圖與學習 dashboard 氛圍}
\end{figure}

\begin{figure}[h]
  \centering
  \includegraphics[width=0.78\linewidth]{docs/img_refs/ECD4C7E2-4BDE-43C0-9A7A-E3F60732EA7A.PNG}
  \caption{Linquest 設計參考圖之二：領地地圖與資訊卡配置方向}
\end{figure}

\subsection{導覽結構}

應用程式採角色導向導覽。學生登入後主要看到 home、roadmap、territory、me 等入口；教師登入後則以 console 為中心，管理班級、活動與結算資料。這樣的設計避免學生看到不必要的管理功能，也讓教師能在 web 版快速完成建立班級、匯入題庫與發布活動。

路由結構使用 Expo Router 的資料夾式 routing。認證相關頁面集中在 \texttt{app/(auth)}，登入後頁面集中在 \texttt{app/(app)}，而 battle room 放在 \texttt{app/battle/[battleId].tsx}，以便從特殊格邀請流程直接進入。

\subsection{色彩方案與設計選擇}

Linquest 的視覺方向是「古典制圖 × 植物圖鑑 × 探險日誌」。專案避免使用高飽和卡通色，而是採用羊皮紙底色、橄欖綠、鼠尾草綠、暖灰與暗陶土紅等低飽和色系。此設計選擇有三個理由：

\begin{enumerate}
  \item 呼應 territory 模式中的地圖、拓荒、佔領與路線感。
  \item 與常見語言學習 app 的明亮卡通風格區隔，建立較沉穩的學習工具感。
  \item 讓教師後台與學生端能共用同一套視覺語言，避免一邊過度遊戲化、一邊過度工具化。
\end{enumerate}

字體方向上，英文標題可使用 Cormorant Garamond 這類古典襯線字，內文與中文則使用清楚易讀的 sans 字體。地圖上的六邊形格會以半透明色塊疊在底圖質感上，倍率格與特殊格則使用細線、金邊或旗幟圖示標示，使學生能直覺分辨策略價值。

\section{技術實作成果}

\subsection{已開發功能}

期中報告中「技術實作進度」仍接近暫無；期末階段已完成多個主流程。以下整理目前專案中的實作狀態。

\begin{longtable}{P{4cm}P{5cm}P{5.5cm}}
\toprule
模組 & 目前狀態 & 對應程式位置 \\
\midrule
認證與角色 & 已有登入、註冊、忘記密碼、session provider 與 role 導向。 & \path{app/(auth)/}、\path{lib/auth/}、\path{lib/ui/session/} \\
\midrule
roadmap & 已有關卡路線、stage 作答、結果頁、Google Sheet 題庫讀取、進度更新與解鎖邏輯。 & \path{app/(app)/roadmap/}、\path{lib/roadmap/} \\
\midrule
territory & 已有地圖產生、格子狀態、佔領規則、refresh、tax、settlement、leaderboard 與 UI 地圖投影。 & \path{app/(app)/territory/}、\path{lib/territory/}、\path{lib/territory-ui/} \\
\midrule
realtime battle & 已有 battle room、邀請、presence、submit answer、watchdog 與 opponent guard。 & \path{app/battle/[battleId].tsx}、\path{lib/realtime-battle/} \\
\midrule
teacher console & 已有班級建立、班級詳情、roster、活動建立、CSV preview、發布、活動詳情、刪除與結算輔助邏輯。 & \path{app/(app)/console/}、\path{lib/teacher-console/}、\path{lib/teacher-console-ui/} \\
\midrule
資料庫與 RLS & 已有 Supabase migrations，涵蓋 auth users、classes、question bank、territory、battle、teacher console、roadmap 與 custom bank。 & \path{supabase/migrations/} \\
\midrule
測試 & 已有大量 unit tests 與 integration tests。integration tests 需獨立 hosted Supabase test project。 & \path{tests/unit/}、\path{tests/integration/} \\
\bottomrule
\end{longtable}

\subsection{程式架構}

專案以 \texttt{app/}、\texttt{lib/}、\texttt{supabase/}、\texttt{tests/} 四個區域分工。

\begin{itemize}
  \item \texttt{app/}：Expo Router 路由與畫面入口。
  \item \texttt{lib/}：服務封裝、領域邏輯、共用 UI 元件與純函式。
  \item \texttt{supabase/migrations/}：資料表、RPC、RLS 與 schema 變更。
  \item \texttt{tests/}：unit tests 與 integration tests。
\end{itemize}

此架構的優點是將商業規則盡量放在 \texttt{lib/} 與 Supabase RPC 中，而不是直接寫在畫面檔案裡。畫面負責呈現狀態與呼叫服務，領地規則、roadmap 解鎖、battle room 狀態、teacher console helper 則各自放在對應模組中，方便測試與維護。

\subsection{資料庫與後端設計}

後端使用 Supabase 作為單一資料來源。使用者、班級、活動、題庫、地圖、作答紀錄與對戰資料都存於 Postgres。RLS 用來限制學生只能讀取自己所屬班級與活動需要的資料；教師則能管理自己建立的班級與活動。

roadmap 與 territory 的題庫設計已分離。roadmap 使用 Google Sheets 作為官方題庫來源，適合維護固定 level；territory 活動則使用教師自訂 CSV 題庫，活動建立時寫入 custom bank，並與活動生命週期綁定。這樣可以避免課堂活動資料污染個人 roadmap 題庫，也讓教師每次活動都能針對本週教材建立內容。

\subsection{測試與品質控管}

專案已配置 Jest 測試。unit tests 涵蓋答題引擎、roadmap、territory pure logic、teacher console helpers、UI components 與 service wrappers。integration tests 涵蓋 territory capture、refresh、tax、settlement、realtime battle、teacher console draft/publish/delete/analytics 等流程。

目前整合測試需要獨立 hosted Supabase test project，因此不應直接指向正式或開發主專案硬跑。這是實務上的限制，但同時也代表專案已經把高風險流程拆成可測試案例，只需要補上隔離測試環境即可完整執行。

\section{專案時程與完成狀態}

\begin{longtable}{P{3.6cm}P{2.7cm}P{3.0cm}P{5.4cm}}
\toprule
任務 & 期中狀態 & 期末狀態 & 說明 \\
\midrule
市場研究 & 已完成 & 已修正與延伸 & 由 Quizlet、Duolingo 比較延伸到教師活動與自訂題庫差異化。 \\
\midrule
UI 設計 & 進行中 & 已建立設計語言並部分落地 & 已有 design language、路由畫面與共用 UI components，仍需視覺 polish。 \\
\midrule
基本應用程式功能 & 尚未完成 & 主流程已完成雛形 & auth、roadmap、territory、battle、teacher console 均有程式實作。 \\
\midrule
資料庫設計 & 初步規劃 & 已有 migrations 與 RLS/RPC & Supabase schema 已涵蓋活動、地圖、題庫、作答、對戰與教師 console。 \\
\midrule
測試與改進 & 尚未開始 & 已建立 unit/integration tests & unit tests 可本機執行；integration tests 需獨立 hosted test project。 \\
\midrule
期末整理 & 未開始 & 已完成報告與 LaTeX 編譯流程 & 本報告以 \texttt{tectonic} 編譯為 PDF。 \\
\bottomrule
\end{longtable}

\section{挑戰與解決方案}

\subsection{遭遇的困難}

\subsubsection{學習性與遊戲性的平衡}

若遊戲規則太弱，學生只是換介面刷題；若遊戲規則太複雜，學生會忽略學習本身。領地、倍率、特殊格、財政、排行榜等元素都需要服務學習目標，而不是只增加操作負擔。

\subsubsection{分數膨脹與公平性}

高價值格若能被兩隊短時間內反覆佔領，就可能出現刷分。特殊格如果太容易被固定強者壟斷，也可能造成弱組失去參與感。

\subsubsection{同步與非同步活動混合}

課堂活動不一定要求全班同時上線，但特殊格 1 對 1 又需要雙方同時在線。這使系統必須同時支援跨時間地圖推進與短時間即時對戰。

\subsubsection{題庫來源分離}

roadmap 題庫與教師活動題庫需求不同。roadmap 需要穩定 level 與進度；教師活動需要快速匯入當次教材。如果共用同一資料模型，會使保留策略與權限控管變複雜。

\subsection{解決方案與調整}

\begin{enumerate}
  \item 地圖格統一大小，改以一般格、倍率格、特殊格區分策略價值，降低規則複雜度。
  \item 導入 cooldown、protected\_until 與 refresh wave 機制，避免高價值格被無限刷分。
  \item 將主要活動設計為非同步，僅特殊格 battle 採 Supabase Realtime，降低整體同步壓力。
  \item roadmap 與 territory 題庫完全解耦：roadmap 讀 Google Sheets，teacher activity 使用 custom bank。
  \item 將 territory/battle attempts 視為活動生命週期資料，roadmap attempts 則保留為未來 SRS 的原料。
\end{enumerate}

\section{教學活動整合與應用}

\subsection{應用程式與學習活動的結合}

Linquest 適合被設計成一週或一單元的課堂活動。教師可以在週初建立活動，貼上本週單字 CSV，設定活動截止時間與小組數。學生在課堂或課後進入活動，透過答題推進團隊領地。教師在活動期間可觀察排行榜與常錯題，並在下一堂課針對弱點補強。

這種方式可支援三種教學活動。

\begin{itemize}
  \item 課堂合作學習：學生分組後共同為地圖推進負責，個人作答會轉化為團隊成果。
  \item 課後練習：教師可將活動期限設定為數天，讓學生在課後持續參與。
  \item 評量與回饋：教師可依活動結算、常錯題與平均正確率，判斷哪些詞彙需要再次教學。
\end{itemize}

\subsection{使用情境模擬}

\subsubsection{情境一：段考前單字複習}

英文老師在段考前一週將考試範圍單字整理為 CSV，建立一輪 Linquest 活動。班上學生自動分為五組，每組在同一張六邊形地圖上拓荒。一般格代表基礎單字，倍率格代表較難或較常錯的單字，特殊格則需要學生找其他組同學對戰。活動結束時，老師依排行榜給予小組獎勵，並用常錯題統計安排考前複習。

\subsubsection{情境二：課後自學與課堂銜接}

學生平時使用 roadmap 模式完成個人關卡，系統記錄答題結果與反應時間。教師在課堂活動中使用另一份自訂題庫進行 territory 活動。兩種模式分別服務「個人練習」與「班級互動」，但都使用四選一答題與單字理解作為核心學習行為。這可讓學生在課外維持練習，也在課堂中獲得社群動機。

\subsection{目標學習成果}

\begin{longtable}{P{4cm}P{10.5cm}}
\toprule
學習成果 & 評量方式 \\
\midrule
知識點掌握 & 透過四選一題目、錯題追問與常錯題統計，觀察學生是否能辨識單字中文與英文對應。 \\
\midrule
反應速度與提取能力 & 記錄 response\_ms，判斷學生是否能快速提取答案，作為未來 SRS 與熟練度估計基礎。 \\
\midrule
合作與互動 & 透過小組領地、特殊格 1 對 1 對戰與排行榜，觀察學生是否願意為團隊參與練習。 \\
\midrule
學習興趣提升 & 以地圖佔領、團隊分數與活動結算作為外在動機，降低重複背誦的枯燥感。 \\
\bottomrule
\end{longtable}

\section{原型展示與執行方式}

專案目前可透過 Expo web 或行動模擬器執行。基本啟動方式如下：

\begin{verbatim}
npm install
cp .env.example .env
npm run web
\end{verbatim}

\texttt{.env} 需填入 Supabase URL、anon key、service role key 與 Google Sheet share URL。若只檢查 unit tests，可執行：

\begin{verbatim}
npm test
\end{verbatim}

若要執行 integration tests，需準備獨立 hosted Supabase test project，避免測試資料影響正式開發資料庫。

\section{結論}

Linquest 從期中的企劃構想，已推進到具備主要程式架構與核心流程的期末原型。系統目前已能清楚呈現產品的兩條主線：一是學生個人 roadmap 自學，二是教師帶班的 territory 領地對戰。前者處理平時練習與關卡進度，後者把課堂題庫轉化為團隊競爭、地圖推進與互動對戰。

本專案最大的特色，是將語言學習中的單字測驗轉化為可被班級共同參與的活動。學生的每一次作答都不只是個人成績，而會影響團隊領地與排行榜；教師也不只是發派題目，而能建立活動、控制題庫、觀察結果並進行回饋。

期末階段仍有可改進之處，包括前端視覺 polish、完整 SRS 演算法、更多題型、正式測試環境與更完整的教師分析儀表板。不過，就期末展示而言，Linquest 已完成從概念、資料模型、後端規則、前端路由到測試結構的主要落地，具備進一步擴充為完整教學產品的基礎。

\section*{參考資料}
\addcontentsline{toc}{section}{參考資料}

\begin{itemize}
  \item \href{https://quizlet.com/en-gb/features/study-modes}{Quizlet Study Modes}
  \item \href{https://www.duolingo.com/learn}{Duolingo Learn}
  \item \href{https://blog.duolingo.com/how-duolingo-streak-builds-habit/}{Duolingo Blog: streak habit design}
  \item \href{https://docs.expo.dev/}{Expo documentation}
  \item \href{https://supabase.com/docs}{Supabase documentation}
  \item 專案內部文件：\path{README.md}
\end{itemize}

\end{document}
